\subsection{由完备赋范李代数定义的群芽}

设 \(H = \sum_{r,s \geqslant 0} H_{r,s} \in \hat{L}(U, V)\) 为豪斯多夫级数（§ 6，第 4 号，定义 1）。我们将证明相应的形式幂级数

\[
\tilde {\mathrm{H}} = \sum_ {r, s \geqslant 0} \tilde {\mathrm{H}} _ {r, s} \in \hat {\mathrm{P}} (\mathfrak {g} \times \mathfrak {g}, \mathfrak {g})\tag{3}
\]

是收敛的（《微分流形与解析流形》R，3.1.1）。

引入下列形式幂级数 \(\eta\in\mathbf{Q}[[U,V]]\)

\begin{enumerate}
\def\labelenumi{(\arabic{enumi})}
\setcounter{enumi}{3}
\tightlist
\item
\item
\end{enumerate}

\[
\begin{array}{l}\eta (\mathrm{U},\mathrm{V}) = -\log (2 - \exp (\mathrm{U} + \mathrm{V}))\\ \qquad = \sum_{m\geqslant 1}\frac{1}{m} (\exp (\mathrm{U} + \mathrm{V}) - 1)^{m}\\ \qquad = \sum_{m\geqslant 1}\frac{1}{m}\sum_{\substack{r_{1},\ldots ,r_{m}\\ s_{1},\ldots ,s_{m}\\ r_{i} + s_{i}\geqslant 1}}\frac{\mathrm{U}^{r_{1}}}{r_{1}!}\frac{\mathrm{V}^{s_{1}}}{s_{1}!}\frac{\mathrm{U}^{r_{2}}}{r_{2}!}\dots \frac{\mathrm{V}^{s_{m}}}{s_{m}!}. \end{array}\tag{6}
\]

因此

\[
\eta (\mathrm{U}, \mathrm{V}) = \sum_ {r, s \geqslant 0} \eta_ {r, s} \mathrm{U} ^ {r} \mathrm{V} ^ {s},\tag{7}
\]

其中

\[
\eta_{r,s} = \sum_{m\geqslant 1}\frac{1}{m}\sum_{\substack{r_{1} + \dots +r_{m} = r\\ s_{1} + \dots +s_{m} = s\\ r_{i} + s_{i}\geqslant 1}}\frac{1}{r_{1}!\ldots r_{m}!s_{1}!\ldots s_{m}!}.\tag{8}
\]

现在令 u 和 v 为满足 \(u + v < \log 2\) 的两个正实数；于是 \(0 \leqslant \exp(u + v) - 1 < 1\)；在 (5) 和 (6) 中分别以 u 代替 U、以 v 代替 V 所得到的级数收敛，且上述计算推出

\[
\sum_ {r, s \geqslant 0} \eta_ {r, s} u ^ {r} v ^ {s} = - \log (2 - \exp (u + v)) <   + \infty .\tag{9}
\]

令 \(r, s \geqslant 0\)，并令 \(\| \tilde{\mathrm{H}}_{r,s} \|\) 为连续多项式 \(\tilde{\mathrm{H}}_{r,s}\) 的范数（《微分流形与解析流形》R，附录，第 2 号）。

引理 1。

\[
\| \widetilde {\mathrm{H}} _ {r, s} \| \leqslant \mathrm{M} ^ {r + s - 1} \eta_ {r, s}.
\]

令 \(r_i, s_i\) 属于 \(\mathbf{N}\)，其中 \(1 \leqslant i \leqslant m\)，且 \(s_m = 1\)；记 \(r = \sum_{i} r_i, s = \sum_{i} s_i\)，并考虑 \(L(\{U, V\})\) 中的下列元素：

\[
Z = \left(\left(\sum_ {i = 1} ^ {m - 1} (\operatorname{ad} U) ^ {r _ {i}} (\operatorname{ad} V) ^ {s _ {i}}\right) (\operatorname{ad} U) ^ {r _ {m}}\right) (V).
\]

于是 \(\tilde{Z} = f \circ p\)，其中 \(f\) 是如下的 \((r + s)\)-线性映射，从 \(\mathfrak{g}^{r + s}\) 到 \(\mathfrak{g}\)：

\[
\begin{array}{l} \left(x _ {1}, \dots , x _ {r}, y _ {1}, \dots , y _ {s}\right) \mapsto \\ \left(\operatorname{ad} \left(x _ {1}\right) \circ \dots \circ \operatorname{ad} \left(x _ {r _ {1}}\right) \circ \operatorname{ad} \left(y _ {1}\right) \circ \dots \circ \operatorname{ad} \left(y _ {s _ {1}}\right) \circ \operatorname{ad} \left(x _ {r _ {1} + 1}\right) \circ \dots \circ \operatorname{ad} \left(x _ {r}\right)\right) \left(y _ {s}\right) \end{array}
\]

其中 \(p\) 是从 \(\mathfrak{g}^2\) 到 \(\mathfrak{g}^{r+s}\) 的如下映射：

\[
(x, y) \mapsto \underbrace {(x , \dots , x} _ {r}, \underbrace {y , \dots , y)} _ {s};
\]

因此 \(\| \tilde{Z}\| \leqslant \| f\| \leqslant M^{r + s - 1}\)（《微分流形与解析流形》R，附录）。将这些不等式应用于 §6 第 4 号公式 (9) 右端的各项，得到：

\[
\| (\mathrm{H} _ {r, s} ^ {\prime}) ^ {\sim} \| \leqslant \frac {\mathrm{M} ^ {r + s - 1}}{r + s} \sum_ {m \geqslant 1} \frac {1}{m} \sum_ {\substack {r _ {1} + \dots + r _ {m} = r \\ s _ {1} + \dots + s _ {m - 1} = s - 1 \\ r _ {1} + s _ {1} \geqslant 1, \dots , r _ {m - 1} + s _ {m - 1} \geqslant 1}} \frac {1}{r _ {1} ! \dots r _ {m} ! s _ {1} ! \dots s _ {m - 1} !}.\tag{10}
\]

类似的论证给出

\[
\begin{array}{l}\| (\mathbf{H}_{r,s}^{\prime \prime})\sim \| \\ \leqslant \frac{\mathbf{M}^{r + s - 1}}{r + s}\sum_{m\geqslant 1}\frac{1}{m}\sum_{\substack{r_{1} + \dots +r_{m - 1} = r - 1\\ s_{1} + \dots +s_{m - 1} = s\\ r_{1} + s_{1}\geqslant 1,\ldots ,r_{m - 1} + s_{m - 1}\geqslant 1}}\frac{1}{r_{1}!\ldots r_{m - 1}!s_{1}!\ldots s_{m - 1}!} \end{array}\tag{11}
\]

从而由 (8)

\[
\left\| \tilde {\mathrm{H}} _ {r, s} \right\| \leqslant   \eta_ {r, s} \frac {\mathrm{M} ^ {r + s - 1}}{r + s} \leqslant \eta_ {r, s} \mathrm{M} ^ {r + s - 1},
\]

这就证明了引理。

命题 1。形式幂级数 \(\tilde{\mathbf{H}}\) 是收敛级数（《微分流形与解析流形》R，3.1.1）；其绝对收敛域（《微分流形与解析流形》R，3.1.4）包含开集

\[
\Omega = \left\{(x, y) \in \mathfrak {g} \times \mathfrak {g} | \| x \| + \| y \| <   \frac {1}{M} \log 2 \right\}.
\]

令 \(u, v\) 为满足 \(> 0\) 的两个 \(u + v < \frac{1}{M} \log 2\) 的实数；则由引理 1

\[
\begin{array}{r l} \mathrm{M} \sum_ {r, s \geqslant 0} \| \tilde {\mathrm{H}} _ {r, s} \| u ^ {r} v ^ {s} \\ & \leqslant \sum_ {r, s \geqslant 0} \eta_ {r, s} \mathrm{M} ^ {r + s} u ^ {r} v ^ {s} = - \log (2 - \exp \mathrm{M} (u + v)) <   + \infty \end{array} \tag {12}
\]

由 (9) 得到。

令 \(h: \Omega \to \mathfrak{g}\) 表示由 \(\tilde{\mathrm{H}}\) 定义的解析函数（《微分流形与解析流形》R，3.2.9），即由公式

\[
h (x, y) = \sum_ {r, s \geqslant 0} \tilde {\mathrm{H}} _ {r, s} (x, y) = \sum_ {r, s \geqslant 0} \mathrm{H} _ {r, s} (x, y) \quad \text { for } (x, y) \in \Omega .\tag{13}
\]

此函数称为关于 M 的 g 的豪斯多夫函数（若不会引起混淆，也简称 g 的豪斯多夫函数）。注意，若 \(\mathrm{H}_{r,s}(\mathrm{U}, -\mathrm{U}) = 0\)，则 \(r + s \geqslant 2\)，从而

\[
h (x, - x) = 0 \quad \text { for } \| x \| <   \frac {1}{2 M} \log 2.\tag{14}
\]

同理

\[
h (0, x) = h (x, 0) = x \quad \text { for } \| x \| <   \frac {1}{\mathrm{M}} \log 2.\tag{15}
\]

命题 2。令

\[
\Omega^ {\prime} = \left\{(x, y, z) \in \mathfrak {g} \times \mathfrak {g} \times \mathfrak {g} | \| x \| + \| y \| + \| z \| <   \frac {1}{M} \log \frac {3}{2} \right\}.
\]

若 \((x,y,z)\in \Omega^{\prime}\)，则

\[
(x, y) \in \Omega , \quad (h (x, y), z) \in \Omega , \quad (y, z) \in \Omega , \quad (x, h (y, z)) \in \Omega\tag{16}
\]

并且

\[
h (h (x, y), z) = h (x, h (y, z)).\tag{17}
\]

令 \((x,y,z)\in \Omega^{\prime}\)；显然 \((x,y)\in \Omega\) 且 \((y,z)\in \Omega\)。此外：

\[
\| h (x, y) \| \leqslant \sum_ {r, s} \| \tilde {\mathrm{H}} _ {r, s} \| \| x \| ^ {r} \| y \| ^ {s},
\]

因而由 (13)

\[
\| h (x, y) \| \leqslant - \frac {1}{M} \log (2 - \exp M (\| x \| + \| y \|)).
\]

现在 \(\mathrm{M}(\|x\| + \|y\|) < \log \frac{3}{2} - \mathrm{M}\|z\|\)；记 \(u = \exp(\mathrm{M}\|z\|)\)；则 \(1 \leqslant u \leqslant \frac{3}{2}\)，并且

\[
\begin{array}{r l} \mathbf {M} (\| h (x, y) \| + \| z \|) <   - \log (2 - \exp (\log \frac {3}{2} - \mathbf {M} \| z \|)) + \mathbf {M} \| z \| \\ = - \log \left(2 - \frac {3}{2 u}\right) + \log u = \log \frac {2 u ^ {2}}{4 u - 3} \\ = \log \left(2 + \frac {2 (u - 1) (u - 3)}{4 u - 3}\right) \leqslant \log 2. \end{array}
\]

同理可见 \((x, h(y, z)) \in \Omega\)。

现在证明 (17)。在李代数 \(\hat{\mathbf{L}}(\{\mathbf{U},\mathbf{V},\mathbf{W}\})\) 中

\[
\mathrm{H} (\mathrm{H} (\mathrm{U}, \mathrm{V}), \mathrm{W}) = \mathrm{H} (\mathrm{U}, \mathrm{H} (\mathrm{V}, \mathrm{W}))
\]

由 §6 第 5 号的命题 4，并由第 1 号公式 (2)，因此在 \(\hat{\mathbf{P}}(\mathfrak{g} \times \mathfrak{g} \times \mathfrak{g}, \mathfrak{g})\) 中有关系

\[
\tilde {\mathrm{H}} \circ (\tilde {\mathrm{H}} \times \mathrm{Id} _ {\mathfrak {g}}) = \tilde {\mathrm{H}} \circ (\mathrm{Id} _ {\mathfrak {g}} \times \tilde {\mathrm{H}}).
\]

由《微分流形与解析流形》R，3.1.9，存在数 \(\varepsilon > 0\)，使得当 \(\|x\|\)、\(\|y\|\) 和 \(\|z\|\) 都 \(\leqslant \varepsilon\) 时公式 (17) 成立。但是，函数 \((x,y,z)\mapsto h(h(x,y),z)\) 和 \((x,y,z)\mapsto h(x,h(y,z))\) 是定义在 \(\Omega'\) 上、取值于 \(\mathfrak{g}\) 的解析函数（《微分流形与解析流形》R，3.2.7）。由于 \(\Omega'\) 连通且它们在 0 的某个邻域内相等，所以它们处处相等（《微分流形与解析流形》R，3.2.5）。

上述结果蕴含：

令 \(\alpha\) 为满足 \(0 < \alpha \leqslant \frac{1}{3M} \log \frac{3}{2}\) 的实数。令

\[
\mathrm{G} = \{x \in \mathfrak {g} \mid \| x \| <   \alpha \},
\]

\(\Theta = \{(x,y) \in G \times G \mid h(x,y) \in G\}\)，并令 \(m: \Theta \to G\) 为 \(h\) 在 \(\Theta\) 上的限制。则：

\begin{enumerate}
\def\labelenumi{(\arabic{enumi})}
\item
  \(\Theta\) 是 \(\mathbf{G} \times \mathbf{G}\) 中的开集，且 \(m\) 是解析的。
\item
  \(x \in G\) 蕴含 \((0, x) \in \Theta\)、\((x, 0) \in \Theta\) 以及 \(m(0, x) = m(x, 0) = x\)。
\item
  \(x \in \mathbf{G}\) 蕴含 \(-x \in \mathbf{G}, (x, -x) \in \Theta, (-x, x) \in \Theta\)，以及
\end{enumerate}

\[
m (x, - x) = m (- x, x) = 0.
\]

\begin{enumerate}
\def\labelenumi{(\arabic{enumi})}
\setcounter{enumi}{3}
\tightlist
\item
  设 \(x, y, z\) 是 \(G\) 中满足 \((x, y) \in \Theta\)、\((m(x, y), z) \in \Theta\)、\((y, z) \in \Theta\) 以及 \((x, m(y, z)) \in \Theta\) 的元素。则 \(m(m(x, y), z) = m(x, m(y, z))\)。
\end{enumerate}

*换言之（第 III 章，§1），若记 \(-x = \sigma(x)\)，则四元组 \((G, 0, \sigma, m)\) 是 \(K_*\) 上的李群芽。

